\section{Conclusion}
\label{sec:conclusion}

Automated mapping of snow avalanche debris in steep alpine terrain is a vital capability for mountain natural hazard risk reduction and emergency response. In this study, we addressed the long-standing challenge of severe topographic-radar distortion in spaceborne Sentinel-1 C-band SAR change detection by proposing \textbf{TopoRadar-Net}, a physics-guided deep change detection network incorporating multi-scale Radar-Topographic Cross-Attention Blocks (RTCAB) and Geomorphically Bounded Loss (Geo-Loss).

Our key scientific and methodological findings include:
\begin{enumerate}
    \item \textbf{Theoretical Formulation of 2D Convolutional Inadequacy in Mountain SAR:} We formulated (Proposition 2) the geometric limitations of shift-invariant 2D spatial linear convolutions across sloped mountain facets, showing that projecting terrain into a continuous 4D aspect-look manifold ($\mathcal{M}$) provides continuous coordinates to condition feature representations without angular branch-cuts.
    \item \textbf{Physics-Guided Cross-Attention Sensitivity:} TopoRadar-Net conditions SAR change representations on directional aspect-look alignment ($\theta_{\text{align}}$) and local incidence angle ($\psi$). Component-removal and stratum analyses quantify sensitivity to these inputs, but do not isolate a causal internal suppression or amplification mechanism.
    \item \textbf{Geomorphic Regularization Sensitivity:} The soft geomorphic boundary penalty regularizes predictions on specified terrain ranges. Removing it produces a $-3.16\%$ paired-cluster F1 difference in Troms{\o} ($p = 0.002$); without valley-specific outcome records or a matched generic regularizer, this result is not attributed to a localized false-positive mechanism.
    \item \textbf{Cross-System Generalization and Stratification Findings:} Across the AvalCD benchmark spanning four global mountain systems, TopoRadar-Net achieves $63.79\% \pm 1.88\%$ pooled macro F1, compared with Attention U-Net ($63.41\%$), SiamUNet-conc ($63.13\%$), ResU-Net ($62.92\%$), and SiamUNet-diff ($54.87\%$). The three-seed gain over SiamUNet-diff is significant ($+8.92\%$, $p = 0.021$), while comparisons with competitive baselines indicate parity ($p = 0.869$ vs Attention U-Net; $p = 0.459$ vs ResU-Net). Seed-42 backslope differences do not remain stable across seeds: the mean differences are $-0.34\% \pm 6.28\%$ in Troms{\o} and $+1.20\% \pm 7.73\%$ in Pamir.
    \item \textbf{Conceptual Decision Support and Quantified Clutter Footprints:} We present an explicitly unvalidated 3-tier conceptual triage framework and quantify a $30.0\%$ ($72.9\text{ ha}$ nominal-grid equivalent, $55.6\text{ ha}$ affine physical reduction) false-alarm difference in the high-altitude Pamirs, alongside $95.8\%$ mean Class D4 completeness in Troms{\o}, a 3.46M-parameter footprint, and $1.80\text{ s}$ model-forward latency.
\end{enumerate}

TopoRadar-Net provides an evaluated all-weather remote-sensing approach for avalanche-debris mapping while explicitly delimiting non-significant competitive-baseline differences, regional failure modes, non-equivariant checkpoint-era augmentation, and the conceptual status of operational triage.
